\textbf{Phần 1. Trắc nghiệm lựa chọn}
\setcounterpageref{ex}{0}

\begin{ex} %[4H2N07]
Đơn vị dùng để đo góc là
\choice
{Xăng-ti-mét}
{Ki-lô-gam}
{\True Độ ($^\circ$)}
{Lít}
\loigiai{
Đơn vị đo góc là độ, kí hiệu là $^\circ$.\\
\textbf{Đáp án của câu hỏi là C}
}
\end{ex}

\begin{ex} %[4H2N07]
Dụng cụ dùng để đo độ lớn của một góc là
\choice
{Thước kẻ thẳng}
{Ê-ke}
{Com-pa}
{\True Thước đo góc (thước đo độ)}
\loigiai{
Thước đo góc dùng để đo độ lớn của góc.\\
\textbf{Đáp án của câu hỏi là D}
}
\end{ex}

\begin{ex} %[4H2H07]
Góc vuông có số đo bằng bao nhiêu độ?
\choice
{$45^\circ$}
{$60^\circ$}
{\True $90^\circ$}
{$180^\circ$}
\loigiai{
Góc vuông có số đo là $90^\circ$.\\
\textbf{Đáp án của câu hỏi là C}
}
\end{ex}

\begin{ex} %[4H2H07]
Góc bẹt có số đo bằng bao nhiêu độ?
\choice
{$90^\circ$}
{\True $180^\circ$}
{$360^\circ$}
{$120^\circ$}
\loigiai{
Góc bẹt có số đo là $180^\circ$.\\
\textbf{Đáp án của câu hỏi là B}
}
\end{ex}

\begin{ex} %[4H2H07]
Góc có số đo $60^\circ$ là
\choice
{Góc vuông}
{Góc tù}
{\True Góc nhọn}
{Góc bẹt}
\loigiai{
Góc $60^\circ < 90^\circ$ nên là góc nhọn.\\
\textbf{Đáp án của câu hỏi là C}
}
\end{ex}

\begin{ex} %[4H2H07]
Góc có số đo $120^\circ$ là
\choice
{Góc vuông}
{\True Góc tù}
{Góc nhọn}
{Góc bẹt}
\loigiai{
Góc $120^\circ$ lớn hơn $90^\circ$ và bé hơn $180^\circ$ nên là góc tù.\\
\textbf{Đáp án của câu hỏi là B}
}
\end{ex}

\begin{ex} %[4H2V07]
Lúc 3 giờ đúng, hai kim giờ và kim phút của đồng hồ tạo thành một góc có số đo là
\choice
{\True $90^\circ$}
{$60^\circ$}
{$120^\circ$}
{$180^\circ$}
\loigiai{
Lúc 3 giờ đúng, hai kim tạo thành góc vuông $90^\circ$.\\
\textbf{Đáp án của câu hỏi là A}
}
\end{ex}

\begin{ex} %[4H2V07]
Lúc 6 giờ đúng, hai kim giờ và kim phút của đồng hồ tạo thành một góc có số đo là
\choice
{$90^\circ$}
{$120^\circ$}
{$150^\circ$}
{\True $180^\circ$}
\loigiai{
Lúc 6 giờ đúng, hai kim thẳng hàng tạo thành góc bẹt $180^\circ$.\\
\textbf{Đáp án của câu hỏi là D}
}
\end{ex}

\begin{ex} %[4H2V07]
Lúc 2 giờ đúng, góc tạo bởi kim giờ và kim phút có số đo là
\choice
{\True $60^\circ$}
{$30^\circ$}
{$90^\circ$}
{$45^\circ$}
\loigiai{
Mỗi khoảng giờ tương ứng $30^\circ$. Lúc 2 giờ góc là $2 \times 30^\circ = 60^\circ$.\\
\textbf{Đáp án của câu hỏi là A}
}
\end{ex}

\begin{ex} %[4H2V07]
Góc bẹt gấp mấy lần góc vuông?
\choice
{$1$ lần}
{\True $2$ lần}
{$3$ lần}
{$4$ lần}
\loigiai{
$180^\circ : 90^\circ = 2$ lần.\\
\textbf{Đáp án của câu hỏi là B}
}
\end{ex}

\begin{ex} %[4H2V07]
Trong một hình tam giác vuông, ngoài 1 góc vuông thì 2 góc còn lại là
\choice
{2 góc tù}
{1 góc nhọn và 1 góc tù}
{2 góc vuông}
{\True 2 góc nhọn}
\loigiai{
Hai góc còn lại của tam giác vuông luôn là hai góc nhọn.\\
\textbf{Đáp án của câu hỏi là D}
}
\end{ex}

\begin{ex} %[4H2C07]
Lúc 4 giờ 30 phút, góc tạo bởi kim giờ và kim phút có số đo là
\choice
{\True $45^\circ$}
{$30^\circ$}
{$60^\circ$}
{$40^\circ$}
\loigiai{
Kim phút chỉ số 6. Kim giờ nằm chính giữa số 4 và số 5 ($4{,}5$).\\
Khoảng cách giữa hai kim: $6 - 4{,}5 = 1{,}5$ khoảng giờ.\\
Số đo góc: $1{,}5 \times 30^\circ = 45^\circ$.\\
\textbf{Đáp án của câu hỏi là A}
}
\end{ex}

\textbf{Phần 2. Trắc nghiệm đúng sai}
\setcounterpageref{ex}{0}

\begin{ex} %[4H2H07]
Xét tính đúng sai về số đo các góc:
\choiceTF
{\True Góc nhọn có số đo bé hơn $90^\circ$.}
{\True Góc vuông có số đo bằng $90^\circ$.}
{Góc tù có số đo lớn hơn $180^\circ$.}
{\True Góc bẹt có số đo bằng $180^\circ$.}
\loigiai{
\begin{itemize}
    \item Góc nhọn $< 90^\circ$. Đúng.
    \item Góc vuông $= 90^\circ$. Đúng.
    \item Góc tù $> 90^\circ$ và $< 180^\circ$ (không thể lớn hơn $180^\circ$). Sai.
    \item Góc bẹt $= 180^\circ$. Đúng.
\end{itemize}
\begin{itemize}
    \item a) \textbf{Đúng}.
    \item b) \textbf{Đúng}.
    \item c) \textbf{Sai}.
    \item d) \textbf{Đúng}.
\end{itemize}
}
\end{ex}

\begin{ex} %[4H2V07]
Xét số đo góc tạo bởi các kim đồng hồ:
\choiceTF
{\True Lúc 1 giờ đúng, góc tạo bởi hai kim là $30^\circ$.}
{\True Lúc 9 giờ đúng, góc nhỏ tạo bởi hai kim là $90^\circ$.}
{Lúc 5 giờ đúng, góc tạo bởi hai kim là góc vuông.}
{\True Lúc 5 giờ đúng, góc nhỏ tạo bởi hai kim là $150^\circ$ (góc tù).}
\loigiai{
\begin{itemize}
    \item Lúc 1 giờ: $1 \times 30^\circ = 30^\circ$. Đúng.
    \item Lúc 9 giờ: góc vuông $90^\circ$. Đúng.
    \item Lúc 5 giờ: $5 \times 30^\circ = 150^\circ$ (góc tù, không phải góc vuông). Ý 3 Sai, Ý 4 Đúng.
\end{itemize}
\begin{itemize}
    \item a) \textbf{Đúng}.
    \item b) \textbf{Đúng}.
    \item c) \textbf{Sai}.
    \item d) \textbf{Đúng}.
\end{itemize}
}
\end{ex}

\textbf{Phần 3. Trả lời ngắn}
\setcounterpageref{ex}{0}

\begin{ex} %[4H2H07]
Tính số đo góc tạo bởi hai kim đồng hồ lúc 4 giờ đúng (theo độ, chỉ ghi số).
\shortans{120}
\loigiai{
\textbf{Đáp số:} $120$\\
$4 \times 30^\circ = 120^\circ$.
}
\end{ex}

\begin{ex} %[4H2H07]
Góc bẹt bằng bao nhiêu độ?
\shortans{180}
\loigiai{
\textbf{Đáp số:} $180$\\
Góc bẹt $= 180^\circ$.
}
\end{ex}

\begin{ex} %[4H2V07]
Một góc có số đo bằng một nửa góc vuông. Tính số đo góc đó (chỉ ghi số).
\shortans{45}
\loigiai{
\textbf{Đáp số:} $45$\\
$90 : 2 = 45^\circ$.
}
\end{ex}

\begin{ex} %[4H2V07]
Tính tổng số đo của một góc vuông và một góc bẹt (theo độ).
\shortans{270}
\loigiai{
\textbf{Đáp số:} $270$\\
$90 + 180 = 270^\circ$.
}
\end{ex}

\begin{ex} %[4H2V07]
Một góc tù có số đo gấp đôi góc $65^\circ$. Hỏi góc đó có số đo bao nhiêu độ?
\shortans{130}
\loigiai{
\textbf{Đáp số:} $130$\\
$65 \times 2 = 130^\circ$.
}
\end{ex}

\begin{ex} %[4H2C07]
Lúc mấy giờ đúng trong khoảng từ 1 giờ đến 6 giờ thì hai kim đồng hồ tạo thành một góc tù có số đo $150^\circ$? (Chỉ ghi số giờ).
\shortans{5}
\loigiai{
\textbf{Đáp số:} $5$\\
$150 : 30 = 5$ giờ đúng.
}
\end{ex}
